\section{The Hausdorff-measure consequence}
\label{sec:consequences}

We now apply the proved Lebesgue-measure theorem to the transformed
radii in the mass transference principle.

\begin{proof}[Proof of Corollary~\ref{cor:hausdorff}]
Put $r_q=\psi(q)/q$. If $r_q\not\to0$, then along an unbounded
sequence of denominators $r_q\ge\delta>0$, so $\psi(q)>1/2$
eventually. Since $\|qx-\gamma\|\le1/2$ for every $x$, in this
case $W_\gamma(\psi)=\mathbb R$.

Suppose now that $r_q\to0$, omit the zero radii, and apply
Theorem~\ref{thm:main} to $\Psi(q)=qf(r_q)$, taking $\Psi(q)=0$
when $r_q=0$. The resulting limsup of open intervals with centres
$c_{a,q}=(a+\gamma)/q$, $a\in\mathbb Z$, and radii $f(r_q)$ has
full Lebesgue measure in every interval.
Retain only centres with $|c_{a,q}|\le q$ and enumerate the resulting
finite rows in increasing $q$. The original and transformed radii
both tend to zero. On every fixed interval $[-M,M]$, every sufficiently
late original or transformed interval meeting it has
$|c_{a,q}|\le M+1<q$, so this restriction changes neither limsup there.

Apply the mass transference principle~\cite[Theorem~2]{BV2006} in
dimension one to the closed balls with these centres and radii $r_q$;
their transformed radii are exactly $f(r_q)$. It gives full
$\mathcal H^f$-measure for the original closed-ball limsup in every
ball. Its difference from the open-ball limsup is contained in the
countable union of all ball endpoints, a set of zero $\mathcal H^f$
measure. This proves the assertion for every bounded interval,
regardless of its endpoint convention; empty and singleton intervals
also have zero measure. The radius convention of~\cite{BV2006}
gives the same assertion as the usual diameter convention: if
$\ell=\lim_{r\downarrow0}f(r)/r<\infty$, both measures are zero or
fixed multiples of length, while if $\ell=\infty$, the
diameter-based measure dominates the radius-based measure, which is
already infinite on the intersection for every nondegenerate interval.
\end{proof}
